\section{Introduction}

Livestock farming is a fundamental cornerstone of global agriculture, food security, and economic stability. As the global population surpasses 8.2 billion in 2025, livestock farming continues to play a vital role in food security, rural livelihoods, and international trade \cite{piyush2025animal}\cite{livestockproduction2026}. Livestock is also a major economic driver. The gross value of global livestock production is projected to increase by 14\% over the next decade, with the livestock sector leading this expansion with a 16\% rise \cite{oecdfao2025outlook}. Global dietary preferences are also shifting toward nutrient-rich animal proteins due to rising incomes and urbanization, especially in middle-income countries in Asia and Latin America \cite{livestockproduction2026}\cite{oecdfao2025outlook}. By 2034, daily per capita consumption of fish and livestock products is expected to increase by 25\% in lower-middle-income nations \cite{oecdfao2025outlook}\cite{livestockproduction2026}.

Swine production is one of the most important sources of animal protein in the world and is heavily regionally concentrated within the larger livestock industry \cite{livestockproduction2026}. Pork is the second most-produced meat after poultry, with an estimated production of 116.27 million metric tons in 2025 \cite{usda2026pork}. China remains the leading producer, with an inventory of approximately 400 million pigs and about 49\% of total global pork production. The European Union (19\%) and the United States (11\%) are the next largest producers \cite{usda2026pork}. Swine farming is also a crucial component of the cultural and economic systems of many developing nations. In the Philippines, for example, the swine industry is valued at around ₱191 billion and accounts for 78\% of the country's total livestock output \cite{livestockproduction2026}\cite{porcinews2026philippine}\cite{lo2025rebooting}. Around 71.1\% of the Philippine hog inventory is managed by backyard smallholder farmers, highlighting the importance of piggery farming to local economic survival \cite{lo2025rebooting}.

The global swine industry is highly susceptible to transboundary diseases, most notably \textbf{African Swine Fever (ASF)}, which has become a major threat to swine-sector growth \cite{livestockproduction2026}. In this context, pig health monitoring is essential for early disease detection, economic sustainability, improved animal welfare, and reduced dependence on antibiotics in agriculture. Although manual inspection of pens, pig behavior, and air quality remains common practice, modern swine farming is steadily shifting toward advanced, continuous monitoring technologies to achieve these goals \cite{beyondtech2025smart}. Since the first confirmed ASF outbreak in the Philippines in July 2019, the disease has wiped out an estimated 6 million pigs, causing the country's swine population to fall to its lowest level in years \cite{livestockproduction2026}\cite{lo2025rebooting}. By 2023, ASF had already caused an estimated ₱100 billion in economic losses \cite{hsu2025perspective}. As local pork production declined sharply, the country became increasingly dependent on meat imports, which reached an estimated 1.64 million metric tons in 2025 \cite{livestockproduction2026}\cite{barro2026philippine}. This dependence on imported pork, combined with supply shortages, contributed to a significant increase in pork prices nationwide.

In response, both public and private sectors have invested in health monitoring and early detection systems to stabilize the industry and limit disease spread. The Department of Agriculture (DA) launched the Bantay ASF sa Barangay (BABay ASF) program to strengthen local monitoring and institutionalize biosecurity practices \cite{hsu2025perspective}\cite{livestockproduction2026}. In addition, DOST-PCAARRD funded the BRIDGES project, which established mobile laboratory units and uses advanced pathogen-detection technologies to provide rapid, on-site disease diagnosis in rural and resource-limited areas \cite{pcaarrd2026}. Manual pig monitoring is highly time-consuming and prone to error, especially when caregivers are overworked. In a case study conducted in a piggery in Tuburan, Cebu, the inability of caretakers to constantly monitor every pig led to the death of 20 nursery pigs in a single month \cite{dingcong2017monitoring}. Meanwhile, commercial operations in regions such as Central Luzon and CALABARZON are increasingly adopting climate-controlled, tunnel-ventilated housing with digital sensors as the sector modernizes. These monitoring conditions have helped reduce pre-weaning mortality from 10\% to 9.43\% while improving productivity and reproductive performance \cite{lo2025rebooting}\cite{livestockproduction2026}.
